%% ========================================================================
%% Brāhmasphuṭasiddhānta of Brahmagupta (628 CE)
%% Part 1, Chunk 6: Pages 701–838 — BILINGUAL EDITION
%% LEFT: Sanskrit in Devanagari  |  RIGHT: English translation with IAST
%% ========================================================================
\documentclass[11pt,a4paper]{article}

%% -------- Core Packages --------
\usepackage{fontspec}
\usepackage{polyglossia}
\usepackage{amsmath,amssymb,amsthm}
\usepackage{geometry}
\usepackage{graphicx}
\usepackage{xcolor}
\usepackage{titlesec}
\usepackage{fancyhdr}
\usepackage{enumitem}
\usepackage{array,booktabs,longtable,multirow}
\usepackage{hyperref}
\usepackage{lettrine}
\usepackage{microtype}
\usepackage{tocloft}
\usepackage{indentfirst}
\usepackage{paracol}
\usepackage{etoolbox}

%% -------- Page Geometry --------
\geometry{
  paperwidth=210mm,paperheight=297mm,
  left=18mm,right=18mm,top=26mm,bottom=26mm,
  headheight=14pt,footskip=30pt
}

%% -------- Language Setup --------
\setmainlanguage{english}
\setotherlanguage{sanskrit}

%% -------- Font Configuration --------
\setmainfont{Linux Libertine O}
\setsansfont{Linux Biolinum O}
\setmonofont{DejaVu Sans Mono}

%% Devanagari (Sanskrit/Hindi) Font — Noto fonts by name
\newfontfamily\devanagarifont[Script=Devanagari]{Noto Serif Devanagari}
\newfontfamily\devanagarifontsf[Script=Devanagari]{Noto Sans Devanagari}
\newfontfamily\devanagarifontbold[Script=Devanagari]{Noto Serif Devanagari Bold}

%% -------- Colors --------
\definecolor{titlecolor}{RGB}{45,55,72}
\definecolor{sectioncolor}{RGB}{55,65,81}
\definecolor{notecolor}{RGB}{120,53,15}
\definecolor{rulecolor}{RGB}{180,160,140}
\definecolor{versecolor}{RGB}{60,40,20}
\definecolor{commentarycolor}{RGB}{40,40,60}
\definecolor{sanskritbg}{RGB}{255,253,245}
\definecolor{englishbg}{RGB}{250,252,255}
\definecolor{dividercolor}{RGB}{160,140,120}

%% -------- Section Formatting --------
\titleformat{\section}
  {\normalfont\Large\bfseries\color{sectioncolor}}
  {\thesection}{1em}{}
\titleformat{\subsection}
  {\normalfont\large\bfseries\color{sectioncolor}}
  {\thesubsection}{1em}{}
\titleformat{\subsubsection}
  {\normalfont\normalsize\bfseries\color{sectioncolor}}
  {\thesubsubsection}{1em}{}

%% -------- Header/Footer --------
\pagestyle{fancy}
\fancyhf{}
\fancyhead[L]{\small\textit{Brāhmasphuṭasiddhānta — Bilingual Edition — Pages 701–838}}
\fancyhead[R]{\small\thepage}
\renewcommand{\headrulewidth}{0.4pt}
\renewcommand{\footrulewidth}{0pt}

%% -------- Custom Commands --------
\newcommand{\longline}{\noindent\rule{\textwidth}{0.4pt}}
\newcommand{\shortline}{\noindent\rule{0.5\textwidth}{0.4pt}\par}
\newcommand{\cnote}[1]{\textcolor{notecolor}{\textit{#1}}}
\newcommand{\dev}[1]{{\devanagarifont #1}}
\newcommand{\devbold}[1]{{\devanagarifontbold #1}}
\newcommand{\devsf}[1]{{\devanagarifontsf #1}}

%% Parallel column environments
\columnratio{0.48}
\setlength{\columnsep}{1.2em}

%% Custom verse environment (no verse.sty needed)
\newenvironment{simpleverse}{%
  \list{}{\itemsep 0pt \parsep 0pt \topsep 4pt \leftmargin 2em}%
  \item\relax\ignorespaces
}{\endlist}

%% Sanskrit verse in left column
\newcommand{\sanskritverse}[1]{%
  \begin{leftcolumn*}
  \begin{simpleverse}\devanagarifont\color{versecolor}#1\end{simpleverse}
  \end{leftcolumn*}
}

%% English translation in right column
\newcommand{\englishverse}[1]{%
  \begin{rightcolumn}
  \begin{simpleverse}#1\end{simpleverse}
  \end{rightcolumn}
}

%% Bilingual verse pair
\newenvironment{bilingualverse}{%
  \begin{paracol}{2}
  \fontsize{10.5}{13}\selectfont
}{%
  \end{paracol}
  \vspace{0.3cm}
}

%% Single-language spanning block
\newcommand{\spanningtext}[1]{%
  \switchcolumn[0]*#1\switchcolumn
}

%% Decorative divider
\newcommand{\divider}{\begin{center}\textcolor{dividercolor}{\rule{0.5\textwidth}{0.5pt}}\end{center}}
\newcommand{\wavydivider}{\begin{center}$\sim\sim\sim\sim\sim\sim\sim\sim\sim\sim\sim\sim\sim\sim\sim\sim\sim\sim\sim\sim$\end{center}}

%% Column labels
\newcommand{\sanskritlabel}{\begin{center}\textcolor{sectioncolor}{\small\textsc{Sanskrit (Devanagari)}}\end{center}}
\newcommand{\englishlabel}{\begin{center}\textcolor{sectioncolor}{\small\textsc{English (with IAST)}}\end{center}}

%% -------- Hyperref Setup --------
\hypersetup{
  colorlinks=true,
  linkcolor=sectioncolor,
  citecolor=sectioncolor,
  urlcolor=sectioncolor,
  pdfencoding=auto,
  pdftitle={Brāhmasphuṭasiddhānta — Bilingual Edition — Pages 701-838},
  pdfauthor={Brahmagupta (628 CE)},
  pdfsubject={Bilingual Sanskrit-English Edition of Brahmasphutasiddhanta Part 1 Final Portion}
}

%% ========================================================================
\newcommand{\vs}{\vspace{0.5em}}
\begin{document}

%% ========================================================================
%% TITLE PAGE
%% ========================================================================
\thispagestyle{empty}
\begin{center}
\vspace*{1.5cm}

{\fontsize{12}{14}\selectfont\dev{श्रीब्रह्मगुप्ताचार्य-विरचित}}\\[0.3cm]
{\fontsize{26}{30}\selectfont\devbold{ब्राह्मस्फुटसिद्धान्तः}}\\[0.4cm]
{\Large\devbold{द्विभाषी संस्करणम्}}\\[0.6cm]
{\large\textit{Brāhmasphuṭasiddhānta — Bilingual Edition}}

\vspace{1.2cm}

\longline

\vspace{0.6cm}

{\LARGE\textbf{Part 1 — Final Portion}}\\[0.4cm]
{\large\textit{Pages 701--838 of the Original Edition}}\\[0.3cm]
{\small Comprising:}\\[0.2cm]
{\normalsize Chapter 23 (Mānādhyāya) · Subject Index ·}\\[0.1cm]
{\normalsize Madhyamādhikāra Commentary}

\vspace{0.6cm}

\longline

\vspace{1cm}

\begin{paracol}{2}
\begin{leftcolumn*}
\begin{center}
{\large\devbold{संस्कृत पाठः}}\\[0.3cm]
{\small\dev{सौरचन्द्रार्क्षसावन-\\मानप्रयोगविधिः}}
\end{center}
\end{leftcolumn*}
\begin{rightcolumn}
\begin{center}
{\large\textbf{Sanskrit Text}}\\[0.3cm]
{\small Solar, Lunar, Stellar, and\\Civil Measures and Application}
\end{center}
\end{rightcolumn}
\end{paracol}

\vspace{1cm}

{\large Author: \textbf{Brahmagupta} (628 CE)}\\[0.3cm]
{\normalsize Editor: Acharya Ramswarup Sharma}\\[0.1cm]
{\small Indian Institute of Astronomical and Sanskrit Research}\\[0.1cm]
{\small New Delhi, First Edition 1966}

\vspace{1.5cm}

{\small Bilingual Edition — Sanskrit (Devanagari) / English (with IAST)}\\[0.2cm]
{\footnotesize Prepared with XeLaTeX · Noto Serif Devanagari · Polyglossia}

\vfill
{\textit{Digital preservation of classical Indian mathematical astronomy}}

\end{center}
\newpage

%% ========================================================================
%% TABLE OF CONTENTS FOR THIS VOLUME
%% ========================================================================
\section*{Contents of this Portion / \dev{अस्यांशस्य विषयसूची}}
\addcontentsline{toc}{section}{Contents of this Portion}

\begin{paracol}{2}
\sanskritlabel
\englishlabel
\switchcolumn[0]*

\begin{enumerate}[label=\dev{\arabic*.},leftmargin=*]
  \item \textbf{\dev{त्रयोविंशो मानाध्यायः}} (पृ.\,७०१--७३९)
  \begin{itemize}[label=\textendash,leftmargin=1em]
    \item \dev{प्रयोगविधिः} — श्लोकाः १--४
    \item \dev{चतुर्विंशः} — श्लोकाः ५--२३
    \item \dev{पदविभाग-विश्लेषणम्}
  \end{itemize}
  \item \textbf{\dev{प्रकाशकस्य विवरणम्}} — पृ.\,७४०--७४९
  \item \textbf{\dev{विषयानुक्रमणिका}} — पृ.\,७५०--७५१
  \item \textbf{\dev{मध्यमाधिकारटीका}} — पृ.\,७५२--८३८
  \begin{itemize}[label=\textendash,leftmargin=1em]
    \item \dev{मध्यमाधिकारस्य विस्तृतसंस्कृत-हिन्दीटीका}
    \item \dev{ग्रहाणां मध्यमगतिविधिः}
  \end{itemize}
\end{enumerate}

\switchcolumn

\begin{enumerate}[label=\arabic*.,leftmargin=*]
  \item \textbf{Chapter 23: Mānādhyāya} (The Measurement Chapter) — pp.\,701--739
  \begin{itemize}[label=\textendash,leftmargin=1em]
    \item Prayogavidhi (Practical Application) — verses 1--4
    \item Caturviṃśa (Twenty-four fold system) — verses 5--23
    \item Detailed word-analysis (Pada-vibhāga) for each verse
  \end{itemize}
  \item \textbf{Title Page and Publisher Information} — pp.\,740--749
  \item \textbf{Viṣayānukramaṇikā} (Subject Index) — pp.\,750--751
  \item \textbf{Madhyamādhikāra} (Mean Computation Chapter) — pp.\,752--838
  \begin{itemize}[label=\textendash,leftmargin=1em]
    \item Extensive Sanskrit–Hindi commentary on mean planetary motion
    \item Table of contents for chapters 1--6 of commentary
  \end{itemize}
\end{enumerate}
\end{paracol}

\divider
\newpage

%% ========================================================================
%% CHAPTER 23: MĀNĀDHYĀYA
%% ========================================================================
\section*{Chapter 23: \dev{त्रयोविंशो मानाध्यायः} — The Measurement Chapter}
\addcontentsline{toc}{section}{Chapter 23: Mānādhyāya (pp.\ 701--739)}
\markboth{Chapter 23: Mānādhyāya (pp.\ 701--739)}{}

\lettrine[lines=3,loversize=0.1]{T}{\textit{he twenty-third chapter}} of the \textit{Brāhmasphuṭasiddhānta} treats of the measurement of time and the practical application of astronomical calculations. This chapter, like others in the work, is composed in the classical \textit{āryā} metre and presents Brahmagupta's rules for calendrical computation, planetary positions, and the reconciliation of different astronomical systems.

The chapter is divided into two principal sections:
\begin{enumerate}[label=(\roman*),leftmargin=2em]
  \item \textbf{Prayogavidhi} — Practical application of the siddhānta
  \item \textbf{Caturviṃśa} — The twenty-four fold division of time
\end{enumerate}

\wavydivider

\subsection*{\dev{प्रयोगविधिः} — Practical Application (Verses 1--4)}

\begin{bilingualverse}
\begin{leftcolumn*}
\dev{सौरैराण्डं मासैस्तिथयश्चान्द्रैंस्तथा सावने दिवसाः ।}\\
\dev{दिनमासाब्दं पमध्यानत दिग्नाकेन्द्रं मानास्म्याम् ॥ १ ॥}
\end{leftcolumn*}
\begin{rightcolumn}
``The solar year, the lunar months and \textit{tithis}, the civil days, days, months, years, the positions of the planets, the \textit{nakṣatras}, and the zodiacal signs---these we call measures (\textit{māna}).''\\[0.3cm]
\hfill\textbf{— Verse 1}
\end{rightcolumn}
\end{bilingualverse}

\textbf{Word Analysis / \dev{पदविभागः}:}

\begin{paracol}{2}
\begin{leftcolumn*}
\dev{(१) सौरैराण्ड} — सौरवत्सरः। \dev{(२) मासैस्तिथयः} — चान्द्रमासाः तिथयः च। \dev{(३) सावनानि दिवसानि} — लौकिकदिनानि। \dev{(४) दिनमासाब्दानि} — अहोरात्रमाससंवत्सराः। \dev{(५) ग्रहनक्षत्रराशयः} — मानानि स्म्याम्॥
\end{leftcolumn*}
\begin{rightcolumn}
(1) \textit{saurairāṇḍa} — the solar year; (2) \textit{māsaistithayaḥ} — lunar months and \textit{tithis}; (3) \textit{sāvanāni divasāni} — civil days; (4) \textit{dinamāsābdāni} — days, months, and years; (5) \textit{grahanakṣatrarāśayaḥ} — planets, nakṣatras, and zodiacal signs---these are called \textit{māna} (measures).
\end{rightcolumn}
\end{paracol}

\vspace{0.4cm}

\begin{bilingualverse}
\begin{leftcolumn*}
\dev{मानान्ति सौरचन्द्रार्क्षां सावनानि ग्रहानयनमेभिः ।}\\
\dev{मानेन पृथक् चतुर्मितिथ्यर्ब्बहारो त्रं लोकस्य ॥ २ ॥}
\end{leftcolumn*}
\begin{rightcolumn}
``The measures of solar, lunar, stellar, and civil reckoning, and the computation of planets by these measures separately---the fourfold manner of computation is in use in the world.''\\[0.3cm]
\hfill\textbf{— Verse 2}
\end{rightcolumn}
\end{bilingualverse}

\begin{paracol}{2}
\begin{leftcolumn*}
\dev{सौरचान्द्रार्क्षसावनानि} — चत्वारि मानानि। \dev{एभिः पृथक् पृथक्} — प्रत्येकम्। \dev{ग्रहानयनम्} — ग्रहाणां गणनम्। \dev{लोकस्य} — संसारे प्रचलितम्॥
\end{leftcolumn*}
\begin{rightcolumn}
\textit{sauracāndrārṣasāvanāni} — four measures; \textit{ebhiḥ pṛthak pṛthak} — separately by each; \textit{grahānanayam} — computation of planets; \textit{lokasya} — current in the world.
\end{rightcolumn}
\end{paracol}

\vspace{0.4cm}

\begin{bilingualverse}
\begin{leftcolumn*}
\dev{युगवर्षं विपुलब्दयनतद्वह्निनिशोद्विद्धिहानयः ।}\\
\dev{सौरास्तिथि करणाऽधिकमासौन राशिपर्व्विक्रियाश्चान्द्रात् ॥ ३ ॥}
\end{leftcolumn*}
\begin{rightcolumn}
``A \textit{yuga-year}, abundant and deficient (days), fire (3), subtraction and addition---these (come) from the solar (system). \textit{Tithi}, \textit{karaṇa}, intercalary month, and the changes of zodiacal signs and \textit{parva} (come) from the lunar (system).''\\[0.3cm]
\hfill\textbf{— Verse 3}
\end{rightcolumn}
\end{bilingualverse}

\begin{paracol}{2}
\begin{leftcolumn*}
\dev{युगवर्षं विपुलाब्दयनानि} — अधिकदिनक्षयदिनानि। \dev{वह्निः} — त्रयः संख्या। \dev{उद्विद्धिहानयः} — क्षेपापचयौ। एते सौरात्॥ \dev{तिथिकरणाधिकमासाः} — चान्द्रात्। \dev{राशिपर्विक्रियाः} — राशिचक्रस्य परिवर्तनानि॥
\end{leftcolumn*}
\begin{rightcolumn}
\textit{yugavarṣaṃ vipulābdayanāni} — intercalary and deficient days; \textit{vahniḥ} — the number three; \textit{udviddhi-hānayaḥ} — addition and subtraction---from the solar system. \textit{tithikaraṇādhikamāsāḥ} --- from the lunar; \textit{rāśiparvikriyāḥ} --- changes of zodiacal signs and lunar half-months.
\end{rightcolumn}
\end{paracol}

\vspace{0.4cm}

\begin{bilingualverse}
\begin{leftcolumn*}
\dev{यसस्य वन प्रमाणं ग्रहगतृपवासं सुतं चिकिर्सतः ।}\\
\dev{सावनं मानांतु शैया प्रायश्चित्त क्रियाश्चान्याः ॥ ४ ॥}
\end{leftcolumn*}
\begin{rightcolumn}
``For one desirous of performing the \textit{pramāṇa} (measurement), the planetary motion, the rising and setting, and the creation (of almanacs), the civil measure (is useful), as well as the bed (auspicious moment), \textit{prāyaścitta} (expiatory rites), and other ceremonies.''\\[0.3cm]
\hfill\textbf{— Verse 4}
\end{rightcolumn}
\end{bilingualverse}

\begin{paracol}{2}
\begin{leftcolumn*}
\dev{प्रमाणं ग्रहगतिक्षयोदयसूतकल्पनम्} — पञ्चाङ्गादिनिर्माणार्थम्। \dev{सावनं मानम्} — लौकिकमानम्। \dev{शय्या} — शयनार्थकालः। \dev{प्रायश्चित्तक्रियाः} — कर्माणि॥
\end{leftcolumn*}
\begin{rightcolumn}
\textit{pramāṇaṃ grahagatikṣayodayasūtakalpanam} --- for the construction of the pañcāṅga (almanac); \textit{sāvanaṃ mānam} --- civil measure; \textit{śayyā} --- auspicious moment for lying down; \textit{prāyaścittakriyāḥ} --- expiatory rites and ceremonies.
\end{rightcolumn}
\end{paracol}

\wavydivider

\subsection*{\dev{चतुर्विंशः} — The Twenty-four Fold System (Verses 5--23)}

\begin{bilingualverse}
\begin{leftcolumn*}
\dev{यस्मात्तत्प्रति पश्तं संज्ञा संज्ञातो विना तस्मात् ।}\\
\dev{लोकप्रसिद्धसांख्यपादीनां दशांकाश्च ॥ ५ ॥}
\end{leftcolumn*}
\begin{rightcolumn}
``Since (computation) is impossible without names (\textit{saṃjñā}) given to each object, the names and numerals well-known in the world (are given).''\\[0.3cm]
\hfill\textbf{— Verse 5}
\end{rightcolumn}
\end{bilingualverse}

\textbf{\dev{टीका} / Commentary:} The commentator explains that Brahmagupta here gives a systematic nomenclature for the various measures of time and computation. The twenty-four items are enumerated in the following verses.

\vspace{0.3cm}

\begin{bilingualverse}
\begin{leftcolumn*}
\dev{युगपद्युगाद्यद् या याम्ययां भास्करस्य वा कन्याम् ।}\\
\dev{राश्यर्द्धसंख्यायां मस्तकया दिनदलाद्वेभ्राम् ॥ ६ ॥}
\end{leftcolumn*}
\begin{rightcolumn}
``A \textit{yugapad}, a \textit{yuga}, beginning from that, or of the sun, of the Yāmya (south), of Bhāskara, or of Kanyā (Virgo), half a sign, the numbers (associated), at the head, from half a day, from two, from Bhrama (revolution).''\\[0.3cm]
\hfill\textbf{— Verse 6}
\end{rightcolumn}
\end{bilingualverse}

\begin{bilingualverse}
\begin{leftcolumn*}
\dev{प्रपमेवकृतः सूर्ये तु पुलिशा रोमक वाराही यवनाद्यैः ।}\\
\dev{यस्मात्तस्मादेकः सिद्धान्तो ग्रथ्य रचनात्मकः ॥ ७ ॥}
\end{leftcolumn*}
\begin{rightcolumn}
``The sun has been measured by Pulisha, Romaka, Vāraha, and Yavana (Greeks). Therefore, from that (measurement) one must construct a proper \textit{siddhānta}.''\\[0.3cm]
\hfill\textbf{— Verse 7}
\end{rightcolumn}
\end{bilingualverse}

\begin{paracol}{2}
\begin{leftcolumn*}
\dev{पुलिशः रोमकः वाराही यवनाः} — पञ्चसिद्धान्तकाराः। \dev{सूर्यस्य प्रमाणम्} — तेजःप्रमाणम्। \dev{रचनात्मकः सिद्धान्तः} — ग्रन्थनिर्माणम्॥
\end{leftcolumn*}
\begin{rightcolumn}
\textit{puliśaḥ romakaḥ vārāhī yavanāḥ} --- the five \textit{siddhānta} authors; \textit{sūryasya pramāṇam} --- measurement of the sun; \textit{racanātmakaḥ siddhāntaḥ} --- construction of a scientific treatise.
\end{rightcolumn}
\end{paracol}

\vspace{0.3cm}

\begin{bilingualverse}
\begin{leftcolumn*}
\dev{यदि भिन्नाः सिद्धान्तभास्कर संस्कान्तयोभि भेदसमाः ।}\\
\dev{सरत्नः पूर्व्वस्यां विपुलत्याकार्द्दयो यस्य ॥ ८ ॥}
\end{leftcolumn*}
\begin{rightcolumn}
``If the \textit{siddhāntas} are different, (yet) Bhāskara (the sun) has the same modifications (\textit{saṃskāra})---the difference (is) abundant and deficient, those having the same (result) for the previous (method).''\\[0.3cm]
\hfill\textbf{— Verse 8}
\end{rightcolumn}
\end{bilingualverse}

\begin{bilingualverse}
\begin{leftcolumn*}
\dev{तत्र पक्षागरिततं मध्यं गत्युत्तरादयः पञ्च ।}\\
\dev{कुडूाकारे बेधः छिन्द्धि चतुर्थां तरे गोलः ॥ ९ ॥}
\end{leftcolumn*}
\begin{rightcolumn}
``In it (the \textit{siddhānta}) the five (methods): \textit{pakṣāgati}, \textit{madhya}, \textit{gati}, \textit{uttara}, etc.; \textit{kuṭūākāra} (transformation), \textit{bheda} (division), \textit{chid} (cutting), \textit{caturtha} (fourth proportional), and the sphere (\textit{gola}).''\\[0.3cm]
\hfill\textbf{— Verse 9}
\end{rightcolumn}
\end{bilingualverse}

\wavydivider

\begin{bilingualverse}
\begin{leftcolumn*}
\dev{भट्ट ब्रह्माचार्येण जिज्ञासुतनयेन गणितगोलविदा ।}\\
\dev{प्रायःश्लेषसंहलं वा स्फुटसिद्धान्तं कृतो ब्रह्मा ॥ १० ॥}
\end{leftcolumn*}
\begin{rightcolumn}
``By Brahmagupta, the son of Jiṣṇu, the knower of mathematics and the sphere, this true \textit{siddhānta} was composed, (which is) Brahma (itself), adorned with the beauty of elegant composition.''\\[0.3cm]
\hfill\textbf{— Verse 10}
\end{rightcolumn}
\end{bilingualverse}

\begin{paracol}{2}
\begin{leftcolumn*}
\dev{भट्ट ब्रह्माचार्यः} — श्रीब्रह्मगुप्तः। \dev{जिज्ञोः तनयः} — जित्स्नुतनयः। \dev{गणितगोलवित्} — गणितखगोलविद्वान्। \dev{स्फुटसिद्धान्तः} — सत्यज्योतिर्ग्रन्थः। \dev{ब्रह्मा} — ब्रह्मस्वरूपः॥
\end{leftcolumn*}
\begin{rightcolumn}
\textit{bhaṭṭa brahmācāryaḥ} --- the venerable Brahmagupta; \textit{jijñoḥ tanayaḥ} --- son of Jiṣṇu; \textit{gaṇitagolavit} --- knower of gaṇita (mathematics) and the celestial sphere; \textit{sphuṭasiddhāntaḥ} --- true astronomical treatise; \textit{brahmā} --- of the nature of Brahman itself.
\end{rightcolumn}
\end{paracol}

\vspace{0.4cm}

\begin{bilingualverse}
\begin{leftcolumn*}
\dev{मग्रहं युति वत्सांकुं विचित्रमलनाद्विव्रहोत्कसमः ।}\\
\dev{शकनिः कर्म्मबाहुत्वान्निर्भुतातो भास्करमग्रहं ॥ ११ ॥}
\end{leftcolumn*}
\begin{rightcolumn}
``(Various) instruments, conjunction, year-count, diverse multiplication, the \textit{utka} (elevation), the \textit{śakani} (depression), on account of the multitude of operations, (all these) from the diminished (part), the solar \textit{maragraham}...''\\[0.3cm]
\hfill\textbf{— Verse 11}
\end{rightcolumn}
\end{bilingualverse}

\begin{bilingualverse}
\begin{leftcolumn*}
\dev{प्राज्येयो नैकल्प्योर्द्धच्छिद्दिने स्थितस्य योग्केस्य ।}\\
\dev{शङ्कुश्ये कथयत्यब्दापि वत्सि सूर्य्यंश्च ॥ १२ ॥}
\end{leftcolumn*}
\begin{rightcolumn}
``The abundant, the one without method, the \textit{ṛddhicchid} (augmented cut), the \textit{yogketsya} of the standing, the \textit{śaṅkuśye} (gnomon-shadow) declares the year and also the solar year...''\\[0.3cm]
\hfill\textbf{— Verse 12}
\end{rightcolumn}
\end{bilingualverse}

\begin{bilingualverse}
\begin{leftcolumn*}
\dev{प्रज मया यन्नोकं गोलादग्रं क्षयार्धसतो त्रं तनु ।}\\
\dev{प्रायःश्लोद्दशोऽयं संज्ञाध्यायचतुर्विंशः ॥ १३ ॥}
\end{leftcolumn*}
\begin{rightcolumn}
``As much as has not been spoken by me, the remainder of the sphere, thin as a sixth part of a \textit{ksaya} (decrease), this (is) the \textit{prāyaśloka} (resumé verse). This (is) the \textit{saṃjñā-adhyāya} (chapter on nomenclature), the twenty-fourth.''\\[0.3cm]
\hfill\textbf{— Verse 13}
\end{rightcolumn}
\end{bilingualverse}

\begin{paracol}{2}
\begin{leftcolumn*}
\dev{मया यन्नोक्तं} — यन्मयोक्तं। \dev{गोलात् अग्रं} — गोलशेषम्। \dev{क्षयार्धसतो षष्ठांशतुल्यं} — अत्यल्पम्। \dev{प्रायःश्लोकः} — सङ्क्षेपश्लोकः। \dev{संज्ञाध्यायः चतुर्विंशः} — नामकल्पनानिर्देशकः अध्यायः॥
\end{leftcolumn*}
\begin{rightcolumn}
\textit{mayā yannoktaṃ} --- what has not been spoken by me; \textit{golāt agraṃ} --- the remainder of the sphere; \textit{kṣayārdhasato ṣaṣṭhāṃśatulyaṃ} --- as minute as one-sixth part of the decrease; \textit{prāyaḥślokaḥ} --- resumé verse; \textit{saṃjñādhyāyaś caturviṃśaḥ} --- the twenty-fourth chapter on nomenclature.
\end{rightcolumn}
\end{paracol}

\vspace{0.3cm}

\textbf{\dev{सूचना} / Note on verse numbering:} The original text numbers this as the 23rd \textit{adhyāya} in the internal count (\dev{इति ब्रह्मगुप्ते मानाध्यायः (२३ अध्याय समाप्तः)}), though the final verse calls it the 24th in another reckoning. The manuscript tradition contains variations in chapter enumeration.

\wavydivider

\begin{bilingualverse}
\begin{leftcolumn*}
\dev{नष्टंघ्नसावनदिनात् सूर्यदिनां त्वसावन दिनानि ।}\\
\dev{यस्मात्तस्मादश्वं दुर्ल्लभगं मन्दबुद्धिनाम् ॥ ५ ॥}
\end{leftcolumn*}
\begin{rightcolumn}
``Multiply the elapsed \textit{sāvana} (civil) days by the solar days; the \textit{asāvana} (non-civil) days---therefore, that attainment is difficult for those of dull intellect.''\\[0.3cm]
\hfill\textbf{— Verse 5 (alt.)}
\end{rightcolumn}
\end{bilingualverse}

\begin{bilingualverse}
\begin{leftcolumn*}
\dev{मानुष्यं पित्रदेवं बाहु यान्यष्टौ वसुर्कालस्य ।}\\
\dev{उत्कर्षि ज्ञानार्थं बाहु स्पष्टं नवममन्यत् ॥ ६ ॥}
\end{leftcolumn*}
\begin{rightcolumn}
``Humanity, the father, the divine, the arm, the eight (\textit{vasu}), the time---the \textit{utkarṣi} (elevation) for the purpose of knowledge, the arm, clearly the ninth (or another).''\\[0.3cm]
\hfill\textbf{— Verse 6 (alt.)}
\end{rightcolumn}
\end{bilingualverse}

\begin{bilingualverse}
\begin{leftcolumn*}
\dev{द्वौ द्वौ राशिं मकराहतवः षट् सूर्यगति वशाऽघोज्याः ।}\\
\dev{शकिरवसन्तं ग्रीष्मा वर्षा शरदौ सहैमन्ताः ॥ ७ ॥}
\end{leftcolumn*}
\begin{rightcolumn}
``Two-two (make) a zodiacal sign; multiplied by Makara (10, i.e.\ 60), these (are) the six (seasons); the solar motion (gives) the seasons: śiśira-vasanta (spring), grīṣma (summer), varṣā (rains), śarad (autumn), and hemanta (winter).''\\[0.3cm]
\hfill\textbf{— Verse 7 (alt.)}
\end{rightcolumn}
\end{bilingualverse}

\begin{paracol}{2}
\begin{leftcolumn*}
\dev{द्वौ द्वौ राशिं} — मिथुनादीनि। \dev{मकराहताः षट्} — दशभिः संख्यया गुणिताः षडृतवः। \dev{सूर्यगतिः} — भानोर्मार्गः। \dev{शिशिर-वसन्त-ग्रीष्म-वर्षा-शरद्-हेमन्ताः} — षडृतवः॥
\end{leftcolumn*}
\begin{rightcolumn}
\textit{dvau dvau rāśiṃ} --- Gemini etc.; \textit{makarāhatāḥ ṣaṭ} --- six seasons multiplied by ten (Makara = 10); \textit{sūryagatiḥ} --- the sun's path; \textit{śiśira-vasanta-grīṣma-varṣā-śarad-hemantāḥ} --- the six seasons.
\end{rightcolumn}
\end{paracol}

\vspace{0.3cm}

\wavydivider

\begin{bilingualverse}
\begin{leftcolumn*}
\dev{मुख्यारः पूजोभक्तः काके व्यासान्तरेण रविकर्णं मुमुच्या ।}\\
\dev{द्वौ, क्षाया दिश्यन्त्वे चन्द्रकर्णेन शेषम् ॥ ८ ॥}
\end{leftcolumn*}
\begin{rightcolumn}
``The principal and the \textit{arha} (worthy), the worshipper and the devotee; by the \textit{vyāsāntara} (diameter difference) for the crow---the sun's hypotenuse is to be released; two, the decrease is to be shown, the remainder by the moon's hypotenuse.''\\[0.3cm]
\hfill\textbf{— Verse 8 (alt.)}
\end{rightcolumn}
\end{bilingualverse}

\begin{bilingualverse}
\begin{leftcolumn*}
\dev{पञ्चदशाहीनयुक्ताः षट् क्रान्तिकालाद्यरसगुणाः ॥ ६४ ॥}
\end{leftcolumn*}
\begin{rightcolumn}
``The six, diminished and increased by fifteen, are the declination-time etc., multiplied by six (\textit{rasa}).''\\[0.3cm]
\hfill\textbf{— Verse 64 (alt.)}
\end{rightcolumn}
\end{bilingualverse}

\begin{bilingualverse}
\begin{leftcolumn*}
\dev{त्रिज्या विपुलछाया घातस्तच्छति विभाजिताक्षाप्तम् ।}\\
\dev{क्षयो नित्यं याम्यो गोलवशाद्विम्बेक्रान्तिः ॥ ६६ ॥}
\end{leftcolumn*}
\begin{rightcolumn}
``The radius and the abundant shadow---their product, divided by that and multiplied by the latitude; the decrease is always in the south (Yāmya); from the sphere, the disc's declination.''\\[0.3cm]
\hfill\textbf{— Verse 66 (alt.)}
\end{rightcolumn}
\end{bilingualverse}

\begin{bilingualverse}
\begin{leftcolumn*}
\dev{क्रान्त्यक्षायुति विभोगाक्षरापदैः शोधिते दिनदलमा ।}\\
\dev{भास्य तिकृतयोः कृतमनुयुतो नयोस्तत्पदे व्यस्ते ॥ ६७ ॥}
\end{leftcolumn*}
\begin{rightcolumn}
``When the sum and difference of declination and latitude are purified by the \textit{akṣarāpada} etc., the half-day; when the square of that is subtracted from the square of the sum---the method when reversed at that position.''\\[0.3cm]
\hfill\textbf{— Verse 67 (alt.)}
\end{rightcolumn}
\end{bilingualverse}

\begin{bilingualverse}
\begin{leftcolumn*}
\dev{पठगुणीतागतं शेषं नाडूच्चै दिवसाङ्कं विभाजितात् त्यात् ।}\\
\dev{दिनदलकर्णगुणाप्ता तया त्रिज्यया फलं कर्णः ॥ ६८ ॥}
\end{leftcolumn*}
\begin{rightcolumn}
``The remainder, multiplied by the prescribed (number) and arrived at, divided by the day-number elevated by the \textit{nāḍī}; that (result) multiplied by the half-day hypotenuse, by that radius the result (is) the hypotenuse.''\\[0.3cm]
\hfill\textbf{— Verse 68 (alt.)}
\end{rightcolumn}
\end{bilingualverse}

\wavydivider

\begin{bilingualverse}
\begin{leftcolumn*}
\dev{इत्युक्तं क्रमेण सांख्याप्रकरणं प्रयोगविधिश्च ।}\\
\dev{व्यासगुणा दीर्घवृत्तं षड़ांशकक्षापयाम् ॥ ९ ॥}
\end{leftcolumn*}
\begin{rightcolumn}
``Thus has been stated in order the section on numerals and the practical application; the diameter multiplied, the long circle, the \textit{ṣaḍāṃśaka} (sixth part), the latitude-decline.''\\[0.3cm]
\hfill\textbf{— Verse 9 (alt.)}
\end{rightcolumn}
\end{bilingualverse}

\begin{bilingualverse}
\begin{leftcolumn*}
\dev{तस्य व्यास शकि करण्डूतद्विजयया गुणालिप्ता ॥ १० ॥}
\end{leftcolumn*}
\begin{rightcolumn}
``Its diameter, śaki ( potency), \textit{karaṇḍū}, that (is) multiplied by the \textit{dvijaya} (twice-conquered), by degrees.''\\[0.3cm]
\hfill\textbf{— Verse 10 (alt.)}
\end{rightcolumn}
\end{bilingualverse}

\begin{bilingualverse}
\begin{leftcolumn*}
\dev{रविकर्णं दूता त्रिजया काके व्यासान्तरं हृता शेषम् ।}\\
\dev{तिजयामुख्यास बधा क्षयि करणं हुतात्मो व्यासः ॥ १० ॥}
\end{leftcolumn*}
\begin{rightcolumn}
``The sun's hypotenuse, the messengers, by the \textit{trijaya} (triple-conquered), by the \textit{vyāsāntara} (diameter difference) divided, the remainder; the \textit{tijayāmukhyāḥ} (chief of the three-conquered), by the decrease, the operation, the \textit{hutātma} (fire-souled = 3), the diameter.''\\[0.3cm]
\hfill\textbf{— Verse 10 (alt. 2)}
\end{rightcolumn}
\end{bilingualverse}

\begin{bilingualverse}
\begin{leftcolumn*}
\dev{व्यासेन दुर्गति बधात् काके व्यासान्तरके मुक्ति बधम् ।}\\
\dev{प्रोक्ष् दुमध्यमुक्तथा तिथिगुणाप्तां तमो व्यासः ॥ ११ ॥}
\end{leftcolumn*}
\begin{rightcolumn}
``From the diameter, the evil, from the product, the crow---the diameter difference, the release (is) the product; the \textit{prokṣa} (sprinkling), the second middle, thus from the \textit{tithi} (lunar day) multiplied and obtained, the darkness, the diameter.''\\[0.3cm]
\hfill\textbf{— Verse 11 (alt.)}
\end{rightcolumn}
\end{bilingualverse}

\begin{bilingualverse}
\begin{leftcolumn*}
\dev{योऽधिकमासावमरात्रसम्भवनः सर्वैर्मानानि ।}\\
\dev{प्रायःश्लोकैर्द्दशाभिर्मानाध्यायश्चतुर्विंशः ॥ १२ ॥}
\end{leftcolumn*}
\begin{rightcolumn}
``The intercalary months, the \textit{avama} (omitted) days, the \textit{sambhavana} (production), all the measures; by ten \textit{prāyaśloka}s (resumé verses), the twenty-fourth Measurement Chapter.''\\[0.3cm]
\hfill\textbf{— Verse 12 (alt.)}
\end{rightcolumn}
\end{bilingualverse}

\vspace{0.5cm}
\begin{center}
\fbox{\dev{\textbf{इति ब्रह्मगुप्ते मानाध्यायः (२३ अध्याय समाप्तः)}}}\\[0.3cm]
\textit{Thus ends the twenty-third chapter on Measurement\\in the work of Brahmagupta.}
\end{center}

\newpage

%% ========================================================================
%% PUBLISHER INFORMATION (Page 745)
%% ========================================================================
\section*{Publisher Information / \dev{प्रकाशकस्य विवरणम्}}
\addcontentsline{toc}{section}{Publisher Information (p.\ 745)}

\begin{paracol}{2}
\sanskritlabel
\englishlabel
\switchcolumn[0]*

\begin{center}
\devbold{प्रकाशकः — Publisher}
\end{center}

\begin{simpleverse}
\dev{इण्डियन इन्स्टीट्यूट ऑफ अस्ट्रानॉमिकल}\\
\dev{एण्ड संस्कृत रिसर्च}\\
\dev{२२३६, गुरुद्वारा रोड, करोल बाग,}\\
\dev{नई दिल्ली-५ (भारत)}
\end{simpleverse}

\switchcolumn

\begin{center}
\textbf{Publisher}
\end{center}

\begin{simpleverse}
Indian Institute of Astronomical\\
and Sanskrit Research\\
2236, Gurudwara Road, Karol Bagh,\\
New Delhi-5 (India)
\end{simpleverse}
\end{paracol}

\vspace{0.5cm}

\begin{paracol}{2}
\begin{leftcolumn*}
\begin{center}
\devbold{सम्पादक मण्डल — Editorial Board}
\end{center}

\begin{simpleverse}
\dev{श्री रामस्वरूप शर्मा}\\
\hspace{2em}\dev{प्रधान सम्पादक, संचालक}

\dev{श्री मुकुन्दमिश्र}\\
\hspace{2em}\dev{ज्योतिषाचार्य}

\dev{श्री विद्यनाथ झा}\\
\hspace{2em}\dev{ज्योतिषाचार्य}

\dev{श्री दयाशंकर दीक्षित}\\
\hspace{2em}\dev{ज्योतिषाचार्य}

\dev{श्री श्रीदत्त शर्मा शास्त्री}\\
\hspace{2em}\dev{एम. ए., एम. श्री. एल.}
\end{simpleverse}
\end{leftcolumn*}

\begin{rightcolumn}
\begin{center}
\textbf{Editorial Board}
\end{center}

\begin{simpleverse}
Sri Ramswarup Sharma\\
\hspace{2em}Principal Editor, Director

Sri Mukundamishra\\
\hspace{2em}Jyotiṣācārya

Sri Vidyanath Jha\\
\hspace{2em}Jyotiṣācārya

Sri Dayashankar Dixit\\
\hspace{2em}Jyotiṣācārya

Sri Shridatt Sharma Shastri\\
\hspace{2em}M.A., M.Shri.L.
\end{simpleverse}
\end{rightcolumn}
\end{paracol}

\vspace{0.5cm}

\begin{paracol}{2}
\begin{leftcolumn*}
\begin{center}
\devbold{प्रथम संस्करण — First Edition: १९६६}

\devbold{मूल्य — Price: ₹ ६०.००}

\devbold{मुद्रक — Printer:}\\
\dev{पं. श्री प्रकाशन एण्ड प्रिन्टर्स}\\
\dev{१२, चमेलियन रोड, दिल्ली}
\end{center}
\end{leftcolumn*}

\begin{rightcolumn}
\begin{center}
\textbf{First Edition: 1966}

\textbf{Price: ₹ 60.00}

\textbf{Printer:}\\
Pandit Shri Prakashan and Printers\\
12, Chameleon Road, Delhi
\end{center}
\end{rightcolumn}
\end{paracol}

\vspace{0.5cm}

\textit{This work was published with financial assistance from the Ministry of Education, Government of India (\dev{भारत सरकार के शिक्षा मन्त्रालय द्वारा प्रदत् अनुदान से प्रकाशित}).}

\newpage

%% ========================================================================
%% TITLE PAGE FOR COMMENTARY (Dvītīya-Bhāga)
%% ========================================================================
\thispagestyle{empty}
\begin{center}
\vspace*{2cm}

{\dev{श्रीब्रह्मगुप्ताचार्य-विरचित}}\\[0.5cm]
{\fontsize{24}{28}\selectfont\devbold{ब्राह्मस्फुटसिद्धान्तः}}\\[0.5cm]
{\large (\textsanskrit{saṃskṛta-hindī-bhāṣayā vāsanā-vijñāna-bhāṣyasya sambhaktaḥ soḍapattrikaḥ})}

\vspace{2cm}

{\LARGE\devbold{द्वितीय-भागः}}\\[0.4cm]
{\Large\textit{Dvītīya-Bhāgaḥ --- Second Part}}

\vspace{0.4cm}
{\large\textit{(Vāsanā-vijñāna Sanskrit--Hindi Commentary)}}

\vspace{2.5cm}

\longline

\vspace{1cm}

{\normalsize\dev{प्रधानसम्पादकः --- Principal Editor}}\\[0.3cm]
{\large\devbold{आचार्यवर पण्डित रामस्वरूप शर्मा}}\\[0.2cm]
{\small (\dev{संचालक-इण्डियन इन्स्टीट्यूट ऑफ अस्ट्रानॉमिकल एण्ड संस्कृत रिसर्च})}

\vspace{1.5cm}

{\normalsize\dev{प्रकाशकः --- Publisher}}\\[0.2cm]
{\dev{इण्डियन इन्स्टीट्यूट ऑफ अस्ट्रानॉमिकल एण्ड संस्कृत रिसर्च}}\\
{\dev{गुरुद्वारा रोड, करोल बाग, न्यू देल्ही-५}}

\end{center}
\newpage

%% ========================================================================
%% SUBJECT INDEX (Viṣayānukramaṇikā)
%% ========================================================================
\section*{\dev{विषयानुक्रमणिका} — Subject Index}
\addcontentsline{toc}{section}{Viṣayānukramaṇikā (Subject Index) (pp.\ 750--751)}
\markboth{Subject Index}{}

\begin{center}
\fbox{\fbox{\dev{\Large\textbf{विषयानुक्रमणिका}}}}
\end{center}

The subject index lists all major topics treated in the commentary (\textit{vāsanā-bhāṣya}) portion of the \textit{Brāhmasphuṭasiddhānta}, with page references to the commentary pages (numbered 1--126 and beyond).

\subsection*{\dev{१. मध्यमाधिकारः} — Mean Computation (pp.\ 1--126)}

\begin{center}
\begin{tabular}{@{}l@{\hspace{1em}}r@{\hspace{2em}}l@{}r@{}}
\toprule
\dhead{Topic / विषयः} & \dhead{Page} & \textbf{English Topic} & \textbf{Page} \\
\midrule
\ditem{मण्डलाचरणम्} & १ & Maṇḍalācaraṇam --- Preliminaries & 1 \\
\ditem{ग्रन्थारम्भप्रयोजनम्} & २ & Granthārambha-prayojanam --- Purpose & 2 \\
\ditem{भचक्रस्य चलनम्} & ३ & Bhacakrasya calanam --- Celestial motion & 3 \\
\ditem{कालप्रवृत्तिः} & ४ & Kāla-pravṛttiḥ --- Commencement of time & 4 \\
\ditem{खण्डस्य विभाग कल्पना} & १४ & Khaṇḍasya vibhāga-kalpanā & 14 \\
\ditem{युगमानम्} & १६ & Yugamānam --- Yuga measurement & 16 \\
\ditem{आर्यभटस्य मतसमीक्षा} & १८ & Āryabhaṭasya mata-samīkṣā & 18 \\
\ditem{मनुप्रमाणानि कल्पप्रमाणम्} & १९ & Manupramāṇāni kalpapramāṇam & 19 \\
\ditem{कल्पविषये आर्यभटमतम्} & २१ & Kalpaviṣaye āryabhaṭamatam & 21 \\
\ditem{रोमकसिद्धान्तमतखण्डनम्} & २३ & Romakasiddhānta-mata-khaṇḍanam & 23 \\
\ditem{ग्रहयुगस्य परिरभाषा} & ३९ & Grahamyugasya paribhāṣā & 39 \\
\ditem{सितरोत्रं शनैश्चरयोः कल्पभगणाः} & ४० & Sitaretro-śanaiścarayoḥ kalpabhagaṇāḥ & 40 \\
\ditem{भच्छ्रमाः कुदिनानि च} & ४२ & Bhacchramāḥ kudināni ca & 42 \\
\ditem{रविवर्षमासराशिमासादयः} & ४४ & Ravivarṣa-māsa-rāśimāsādayaḥ & 44 \\
\ditem{सावननक्षत्रमानवदिनानि} & ४६ & Sāvananakṣatramānavadināni & 46 \\
\ditem{गतकालस्याहर्गणादीनां परिहारः} & ५८ & Gatakālasya āhargaṇādīnāṃ parihāraḥ & 58 \\
\ditem{आर्यभटमतम्} & ५८ & Āryabhaṭamatam --- Āryabhaṭa's view & 58 \\
\ditem{कल्पगताक्षावनाहर्गणाः} & ५९ & Kalpagatākṣāvanāhargaṇāḥ & 59 \\
\ditem{ग्रहमन्दादिमध्यमानां मध्यमानयनम्} & ६५ & Grahamandādi-madhyamānāṃ madhyamānayanam & 65 \\
\ditem{स्वसिद्धान्तनिर्देशनम्} & ७६ & Svasiddhānta-nirdeśanam & 76 \\
\ditem{प्रश्नरात्रौ वारप्रवृत्तिपक्षदर्शनम्} & ७६ & Praśnarātrau vāra-pravṛtti-pakṣadarśanam & 76 \\
\ditem{मध्यमानां देशनियमार्थः} & ७८ & Madhyamānāṃ deśaniyamārthaḥ & 78 \\
\ditem{तदर्थं देशान्तरकर्म्मणा प्रदर्शनम्} & ८२ & Tadarthaṃ deśāntarakarmaṇā pradarśanam & 82 \\
\ditem{वर्षदीर्घ दिनाध्यमदिनादिसाधनम्} & ८६ & Varṣadīrgha dinādhyamadinādi-sādhanam & 86 \\
\ditem{वर्षदीर्घमाससाधनम्} & ८८ & Varṣadīrghamāsa-sādhanam & 88 \\
\ditem{वर्षशानयनं लध्वहर्गणानयनम्} & ९१ & Varṣaśānayanaṃ ladhvahargaṇānayanam & 91 \\
\ditem{रविसितबुधानां कुजगुरुशनिरीश्रोक्ष्णां चानयनम्} & ९६ & Ravisitabudhānāṃ kujaguruśanirīśrekṣṇāṃ cānanayam & 96 \\
\ditem{मध्यमचन्द्रानयनम्} & १०० & Madhyamacandrānayanam & 100 \\
\ditem{वर्षान्तिकाद्धर्गणात् भौमादिग्रहमन्दफलानयनम्} & १०२ & Varṣāntikād dhargaṇāt bhaumādigrahamandaphalānayanam & 102 \\
\bottomrule
\end{tabular}
\end{center}

\newpage

\subsection*{\dev{२. स्फुटाधिकारः} — True Position Computation (pp.\ 160--255)}

\begin{center}
\begin{tabular}{@{}l@{\hspace{1em}}r@{\hspace{2em}}l@{}r@{}}
\toprule
\dhead{Topic / विषयः} & \dhead{Page} & \textbf{English Topic} & \textbf{Page} \\
\midrule
\ditem{स्फुटीकरणस्य प्रयोजनम्} & १६० & Sphuṭīkaraṇasya prayojanam --- Purpose & 160 \\
\ditem{आर्धज्योत्क्रमज्याकथनम्} & १६० & Ārdhjyotkramajyā-kathanam & 160 \\
\ditem{चापज्यानयनम्} & १६९ & Cāpajyānanayam --- Arc-sine computation & 169 \\
\ditem{ज्यातच्छापानयनम्} & १९१ & Jyāta-ccāpānayanam & 191 \\
\ditem{मन्दशीघ्रकेन्द्रयोः केन्द्रभुजकोटिज्ययो व्य परिर्मापा} & १९५ & Mandaśīghrakeśrayoḥ... & 195 \\
\ditem{स्फुटपरिरक्ष्यानयनम्} & १९९ & Sphuṭaparirakṣyānayanam & 199 \\
\ditem{रविचन्द्रयोः स्फुटत्वार्थः} & १९८ & Ravicandrayoḥ sphuṭatvārthaḥ & 198 \\
\ditem{आर्यभटोक्तं स्फुटीकरणम्युत्कम्} & १९८ & Āryabhaṭoktaṃ sphuṭīkaraṇamyutkam & 198 \\
\ditem{रविचन्द्रयोः स्फुटीकरणार्थं मन्दपरिरक्ष्यांशाः} & १९९ & Mandaparirakṣyāṃśāḥ & 199 \\
\ditem{दृष्टे नते स्फुटपरिरक्ष्यानयनम्} & २८४ & Dṛṣṭe nate sphuṭaparirakṣyānayanam & 284 \\
\ditem{मन्दफलस्य धनत्वमृणात्वं च} & २८५ & Mandaphalasya dhanatvamṛṇātvaṃ ca & 285 \\
\ditem{रविचन्द्रयोः स्फुटीकरणो विशेषः} & २८७ & Ravicandrayoḥ sphuṭīkaraṇo viśeṣaḥ & 287 \\
\ditem{ग्रहयोः सूर्याचन्द्रमसोर्न्तकालः} & १९२ & Grahamyoḥ sūryācandramasorntakālaḥ & 192 \\
\bottomrule
\end{tabular}
\end{center}

\subsection*{\dev{द. चन्द्र छायाधिकारः} — Lunar Shadow Chapter (pp.\ 501--548)}

\begin{center}
\begin{tabular}{@{}l@{\hspace{1em}}r@{\hspace{2em}}l@{}r@{}}
\toprule
\dhead{Topic / विषयः} & \dhead{Page} & \textbf{English Topic} & \textbf{Page} \\
\midrule
\ditem{आरम्भप्रयोजनम्} & ५०१ & Ārambha-prayojanam & 501 \\
\ditem{कर्तव्यताकथनम्} & ५०१ & Kartavyatā-kathanam & 501 \\
\ditem{चन्द्रस्य दूयाद्दूयत्वम्} & ५०२ & Cakrasya dūyād-dūyatvam & 502 \\
\ditem{चन्द्रैतकालसाधनम्} & ५०३ & Caitrakāla-sādhanam & 503 \\
\ditem{चन्द्रशङ्क्वानयनम्} & ५०८ & Candraśaṅkvānayanam & 508 \\
\ditem{रविचन्द्रयोः स्फुटशङ्क्वानयनम्} & ५०५ & Ravicandrayoḥ sphuṭaśaṅkvānayanam & 505 \\
\ditem{चन्द्रस्य स्फुटच्छायासाधनप्रकारः} & ५०८ & Sphuṭacchāyā-sādhanaprakāraḥ & 508 \\
\ditem{मध्यच्छायासाधनप्रकारः} & ५०९ & Madhyacchāyā-sādhanaprakāraḥ & 509 \\
\ditem{अध्यायोपसंहारः} & ५११ & Adhyāyopasaṃhāraḥ --- Chapter conclusion & 511 \\
\bottomrule
\end{tabular}
\end{center}

\subsection*{\dev{६. ग्रहयुत्यधिकारः} — Planetary Conjunction Chapter (pp.\ 515--567)}

\begin{center}
\begin{tabular}{@{}l@{\hspace{1em}}r@{\hspace{2em}}l@{}r@{}}
\toprule
\dhead{Topic / विषयः} & \dhead{Page} & \textbf{English Topic} & \textbf{Page} \\
\midrule
\ditem{ग्रहणां मध्यमशरकला मध्यमविम्बकला} & ५१५ & Madhyamaśarakalā & 515 \\
\ditem{ग्रहविम्बकला स्फुटीकरणम्} & ५१७ & Grahavimbakalā sphuṭīkaraṇam & 517 \\
\ditem{युतिकालज्ञानार्थं चालनफलज्ञानार्थंश्र} & ५२१ & Yuti-kāla-jñānārtham & 521 \\
\ditem{चालनफलसंस्कारदारा ग्रहयोः सम्लिप्तीकरणार्थः} & ५२३ & Cālanaphalasaṃskāra-dvārā... & 523 \\
\ditem{स्फुटपातानयनम्} & ५२४ & Sphuṭapātānayanam & 524 \\
\ditem{गणितागतदेवपाताद्यमध्यमसंस्कारसाधनोपायः} & ५२७ & Gaṇitāgata-devapātādi... & 527 \\
\ditem{युतिकाले ग्रहशरसाधनम्} & ५२८ & Yutikāle grahaśara-sādhanam & 528 \\
\ditem{युतिकाले ग्रहयोद्दिष्टोपान्तरः} & ५४४ & Yutikāle grahayoddiṣṭopāntaraḥ & 544 \\
\bottomrule
\end{tabular}
\end{center}

\newpage

%% ========================================================================
%% MADHYAMĀDHIKĀRA COMMENTARY
%% ========================================================================
\section*{\dev{मध्यमाधिकारः} — Mean Computation Chapter}
\addcontentsline{toc}{section}{Madhyamādhikāra Commentary (pp.\ 752--838)}
\markboth{Madhyamādhikāra Commentary}{}

\lettrine[lines=3,loversize=0.1]{T}{\textit{he Madhyamādhikāra}} constitutes the first and foundational chapter of the commentary (\textit{vāsanā-vijñāna-bhāṣya}) on the \textit{Brāhmasphuṭasiddhānta}. It treats the computation of mean positions of planets (\textit{madhyama-graha}), beginning with a discussion of the motion of the celestial sphere (\textit{bhacakra}), the definition of time (\textit{kāla-pravṛtti}), and the various \textit{yuga}-systems employed in Indian astronomy.

\subsection*{Introductory Discussion / \dev{आरम्भकथनम्} (pp.\ 752--760)}

The commentary opens with a discussion of the \textit{bhacakra} (celestial sphere) and its motion:

\begin{bilingualverse}
\begin{leftcolumn*}
\dev{प्रदक्षिणागम् (दक्षिणावस्स्क्ष्मेण चलायमानम्) श्रादो (सर्व्वप्रथमं) ब्रह्मणा}\\
\dev{(जगद्गुरुपादकेन) सृष्टम् (रचितम्) ग्रयदिध्नयादौ स्थितेनक्षत्रादिमततुल्यचपा-}\\
\dev{तादिमित्र साकं भ्रुवयष्ट्याद्यारेण प्रमाणशील ज्योतिष्मत् प्रदायिनां नक्षत्रादीनां}\\
\dev{चक्रं (गोलं) ब्रह्मणा रचितं यथाप्याचार्ययोः कादीनां वचां न कियते तेषां स्वरूपा-}\\
\dev{भावात् किन्तु स्पष्टादिकाले तेषामपि स्थानाऽऽज्ञानरूपसर्जें कृतमिति ॥ १ ॥}
\end{leftcolumn*}
\begin{rightcolumn}
``(The celestial sphere), moving eastward (from left to right), was first created by Brahmā (the world-teacher). The sphere, along with the nakṣatras and other celestial bodies having equal \textit{capādas}, along with the meridian staff etc., the circle of nakṣatras and other luminaries that give light---this sphere was fashioned by Brahmā. Even so, the ācāryas (teachers) do not describe their (the planets') essential nature (\textit{svārūpa}), but have only determined their positions (\textit{sthāna}) for the purpose of computation (\textit{spaṣṭādi}).''\\[0.3cm]
\hfill\textbf{— Commentary v.\ 1}
\end{rightcolumn}
\end{bilingualverse}

\textbf{\dev{हिन्दीटीका} / Hindi Commentary (pp.\ 752--753):}
The commentator explains that the celestial sphere (\textit{bhacakra}) was first created by Brahmā (the creator). The sphere, along with the nakṣatras and other celestial bodies, revolves from west to east. The commentator notes that the \textit{ācāryas} (teachers) do not discuss the nature (\textit{svārūpa}) of these bodies but have only determined their positions (\textit{sthāna}) for the purposes of computation.

\wavydivider

\subsection*{Key Verses from the Siddhānta / \dev{सिद्धान्तमूलश्लोकाः} (pp.\ 760--770)}

The commentary elucidates the following root text (\textit{mūla}) verses:

\begin{bilingualverse}
\begin{leftcolumn*}
\dev{इदानीमनाख्यानतस्य कालस्य प्रवृत्तिमाह ।}\\
\dev{चैनंसितादेर्द्दयाद्व मानोर्द्धनमासबर्षयुगकल्पः ।}\\
\dev{सृष्ट्यादौ लङ्कायां स्मं प्रवृतta दिनेस्कस्य ॥ ४ ॥}
\end{leftcolumn*}
\begin{rightcolumn}
``Now (Brahmagupta) declares the commencement of that time whose beginning is unmentioned. Day, month, year, \textit{yuga}, and \textit{kalpa} began from the sun at Laṅkā at the beginning of creation.''\\[0.3cm]
\hfill\textbf{— Mūla-śloka 4}
\end{rightcolumn}
\end{bilingualverse}

\textbf{Hindi commentary (pp.\ 761--762):}

\begin{paracol}{2}
\begin{leftcolumn*}
\dev{चैत्यादि कल्पप्रतिपद्यात् लंकोपलक्षितांभृद्वेकोर्कोदयादिदिवा प्रवृत्ताः} — चैत्रशुक्लप्रतिपदादिकालात् लङ्कायां सूर्योदयात् दिनादयः प्रवृत्ताः।

\dev{रविदिनम्बरः सृष्ट्यादौ कल्पादौ} — सृष्ट्यारम्भे कल्पस्य आदौ रवेर्दिनसंख्या।
\end{leftcolumn*}
\begin{rightcolumn}
\textit{caityādi kalpa-pratipadyāt laṅkopalakṣitāṃbhṛdv ekodayādi-divā pravṛttāḥ} --- From the first tithi of Caitra (the first month), at Laṅkā, the days etc.\ commenced from the sunrise.

\textit{ravidinambaraḥ sṛṣṭyādau kalpādau} --- The solar day number at the beginning of creation, at the beginning of the \textit{kalpa}.
\end{rightcolumn}
\end{paracol}

\vspace{0.3cm}
The commentator explains the concept of the \textit{kuldiņa} (civil day) in detail. The number of civil days elapsed since the beginning of the current kalpa is 1,58,22,23,78,50,000 according to Brahmagupta's system.

\wavydivider

\subsection*{On Time Computation and Planetary Motion / \dev{कालगणनम् ग्रहचलनं च} (pp.\ 770--790)}

The commentary proceeds with detailed mathematical analysis:

\begin{bilingualverse}
\begin{leftcolumn*}
\dev{प्रथमं युगपदं चलनं द्वितीयं पश्चादिनपर्व्वैः}\\
\dev{तृतीयं तु त्रयोदशभिर्मासैः सावनेन वा चतुर्थम् ॥}
\end{leftcolumn*}
\begin{rightcolumn}
``First, the \textit{yugapada} (yuga-epoch method); second, by \textit{parva} (lunisolar method); third, by thirteen months; or fourth, by \textit{sāvana} (civil reckoning).''\\[0.3cm]
\hfill\textbf{— Time-reckoning verse}
\end{rightcolumn}
\end{bilingualverse}

\textbf{\dev{व्याख्या} / Exposition:} The four types of time-reckoning are discussed:
\begin{enumerate}[label=(\roman*),leftmargin=2em]
  \item \textit{Yugapada} --- the yuga-epoch method
  \item \textit{Parva} --- the lunisolar intercalation method
  \item \textit{Trayodaśa-māsa} --- the thirteen-month intercalary method
  \item \textit{Sāvana} --- civil reckoning based on sunrise
\end{enumerate}
Each method is explained with its mathematical basis and astronomical application.

\wavydivider

\subsection*{Planetary Computations / \dev{ग्रहगणनविधिः} (pp.\ 790--820)}

The commentary provides detailed algorithms for computing mean planetary positions:

\textbf{Mean longitude formula / \dev{मध्यमीभावनसूत्रम्}:}
\begin{equation*}
\text{Mean longitude} = \frac{\text{ahargaṇa} \times \text{daily motion}}{\text{civil days in a kalpa}} \mod 360°
\end{equation*}

Where:
\begin{itemize}[leftmargin=2em]
  \item \textit{ahargaṇa} = elapsed civil days since kalpa-beginning
  \item \textit{daily motion} = mean daily motion of the planet
  \item The result is expressed in degrees of the ecliptic
\end{itemize}

The commentary discusses the \textit{mandaphala} (equation of centre) and \textit{śīghraphala} (\textit{śīghra} equation) with detailed trigonometric formulas using the \textit{jyā} (sine) function.

\begin{paracol}{2}
\begin{leftcolumn*}
\dev{मन्दफलं रविचन्द्रयोः — मन्दकendraस्य ज्या गुणकः, त्रिज्या भाजकः।}

\dev{शीघ्रफलं बुधगुरुशनिभ्यः — शीघ्रकेन्द्रस्य ज्या गुणकः।}
\end{leftcolumn*}
\begin{rightcolumn}
\textit{Mandaphala} for the Sun and Moon: the sine of the \textit{manda-kendra} (anomaly) is the multiplier, the radius (\textit{trijyā}) is the divisor.

\textit{Śīghraphala} for Mercury, Jupiter, and Saturn: the sine of the \textit{śīghra-kendra}.
\end{rightcolumn}
\end{paracol}

\wavydivider

\subsection*{On the Computation of Weekdays / \dev{वारप्रवृत्तिकथनम्} (pp.\ 820--838)}

The final section of this portion treats the computation of weekdays (\textit{vāra-pravṛtti}):

\begin{bilingualverse}
\begin{leftcolumn*}
\dev{वारप्रवृत्तिं मनयो वदन्ति सूर्योदयाद्विष्णुराज्ञाध्याम् ।}\\
\dev{उर्वी तथाऽघोषपरत्र तस्याऽर्चराथं देशान्तरनाडिकाभिः ॥ १३६॥}
\end{leftcolumn*}
\begin{rightcolumn}
``The sages declare the commencement of the week (to be reckoned) from the sunrise (at Laṅkā); on the earth (at other places), it (is determined) by adding or subtracting the \textit{deśāntara-nāḍikās} (longitude difference in time).''\\[0.3cm]
\hfill\textbf{— Verse 136}
\end{rightcolumn}
\end{bilingualverse}

\begin{paracol}{2}
\begin{leftcolumn*}
\dev{लङ्कोदयाम्यसुजातप्रथममपरतः पूर्व्वदेशे च पश्चात्—}\\
\dev{दर्शवोत्थानिमष्टटीभिः सवितः क्षदयतो वासरैव प्रवृत्तिः ॥}\\
\dev{यं या सूर्योदयात् प्राक् चरराकलम्बवैश्चास्मियमियगोले ।}\\
\dev{पश्चात् सौम्यगोले युतिर्वियुतिर्वशाच्छोभयोः स्पष्टकालः इति ॥१३६॥}
\end{leftcolumn*}
\begin{rightcolumn}
``At Laṅkā, the weekday begins with the first \textit{nāḍī} after sunrise; at a place to the west, (it begins) after (sunrise); with the six or eight \textit{nāḍī}s (difference), the weekday is determined when the sun rises. For a place on the eastern or western side of the prime meridian, before or after sunrise (at Laṅkā), by addition or subtraction, the correct time (is found).''\\[0.3cm]
\hfill\textbf{— Expanded v.\ 136}
\end{rightcolumn}
\end{bilingualverse}

\textbf{\dev{हिन्दीटीका} / Hindi Commentary (pp.\ 836--838):}

\begin{paracol}{2}
\begin{leftcolumn*}
\dev{लङ्कायां सूर्योदयात् प्रथमनाड्याम् वारारम्भः। पश्चिमदेशे पश्चात् षडष्टघटीभिः।}

\dev{प्राच्यां देशे सूर्योदयात् प्राक् स्पष्टकालः।}
\end{leftcolumn*}
\begin{rightcolumn}
At Laṅkā, the weekday begins at sunrise in the first \textit{nāḍī}. At a place to the west, after (sunrise) by six or eight \textit{ghaṭī}s.

At a place to the east, the correct time (is found) before sunrise.
\end{rightcolumn}
\end{paracol}

The commentator explains that the \textit{vāra} (weekday) begins at sunrise. By this definition, at a place east of the prime meridian (Laṅkā), the weekday begins after subtracting the longitude difference (in \textit{ghaṭī}s); at a place west of the meridian, (the weekday begins) before sunrise by the same amount (longitude difference in \textit{ghaṭī}s).

In the \textit{Siddhānta-śiromaṇi}, Bhāskarācārya states: ``\textit{arkodayād dvābhyāṃ bhyas tāniḥ prācyāṃ praticyāṃ dinap-}''---(the computation proceeds from sunrise with modifications for eastern and western positions).

\wavydivider

\subsection*{Concluding Verses / \dev{उपसंहारश्लोकाः} (p.\ 838)}

The portion concludes with a verse from the \textit{Manusmṛti}:

\begin{bilingualverse}
\begin{leftcolumn*}
\dev{आसीदिदं तमोभूतमप्रज्ञातमलक्षणम् ।}\\
\dev{अप्रतर्क्यमनाऽऽगम्यं प्रसुप्तमिव सर्व्वतः ॥ इति ॥}
\end{leftcolumn*}
\begin{rightcolumn}
``This (universe) existed as darkness, unknowable, without marks, not to be reasoned about, not to be apprehended, as if asleep everywhere.''\\[0.3cm]
\hfill\textbf{— Manusmṛti 1.5}
\end{rightcolumn}
\end{bilingualverse}

\textbf{\dev{टीकाकारवचनम्} / Commentator's exposition:}

The commentator explains that Manu has stated that at the time of the deluge (\textit{pralaya}), the perception of weekdays and other time measures is not possible. Therefore, the reckoning of weekdays is only possible during the manifest (\textit{sṛṣṭi}) period. Brahmagupta establishes that the weekday begins from sunrise at Laṅkā and is modified by the longitude difference (\textit{deśāntara}) for other locations.

\newpage

%% ========================================================================
%% APPENDIX: STRUCTURAL OVERVIEW OF THE COMMENTARY
%% ========================================================================
\section*{Appendix: Structural Overview of the Commentary / \dev{टीकायाः संरचनासूची}}
\addcontentsline{toc}{section}{Appendix: Structural Overview}

The \textit{Vāsanā-vijñāna} commentary on the \textit{Brāhmasphuṭasiddhānta} is organized into the following major chapters (\textit{adhyāya}s):

\begin{center}
\begin{tabular}{@{}clc@{}}
\toprule
\dhead{No.} & \dhead{Title / \dev{शीर्षकः}} & \textbf{Pages} \\
\midrule
\ditem{१} & \ditem{मध्यमाधिकारः} --- Mean Computation & 1--126 \\
\ditem{२} & \ditem{स्फुटाधिकारः} --- True Position Computation & 160--255 \\
\ditem{३} & \ditem{त्रिप्रश्नाधिकारः} --- Three Questions (Diurnal) & 256--344 \\
\ditem{४} & \ditem{चन्द्रग्रहणाधिकारः} --- Lunar Eclipse & 345--400 \\
\ditem{५} & \ditem{चन्द्रच्छायाधिकारः} --- Lunar Shadow & 501--548 \\
\ditem{६} & \ditem{ग्रहयुत्यधिकारः} --- Planetary Conjunction & 515--567 \\
\ditem{७} & \ditem{नक्षत्रच्छायाधिकारः} --- Stellar Shadow & --- \\
\ditem{८} & \ditem{ग्रहच्छायाधिकारः} --- Planetary Shadow & --- \\
\ditem{९} & \ditem{उदयास्तमयाधिकारः} --- Heliacal Rising/Setting & --- \\
\ditem{१०} & \ditem{पाताधिकारः} --- Nodes and Intersections & --- \\
\ditem{११} & \ditem{गणिताध्यायः} --- Mathematics (Chs.\ 12--18) & --- \\
\ditem{१२} & \ditem{कुण्डकाध्यायः} --- Kuṭṭaka (Pulverizer) & --- \\
\ditem{१३} & \ditem{परिकर्माध्यायः} --- Auxiliary Computations & --- \\
\bottomrule
\end{tabular}
\end{center}

\divider

\vspace{0.5cm}
\begin{center}
\textit{The commentary employs a rich methodology, interweaving:}
\begin{enumerate}[label=(\roman*),leftmargin=2em]
  \item \textit{Mūla-śloka} --- The root text verses of Brahmagupta
  \item \textit{Anvaya} --- Word-order rearrangement for grammatical clarity
  \item \textit{Bhāṣya} --- Sanskrit explanatory commentary
  \item \textit{Hindi Ṭīkā} --- Hindi translation and explanation
  \item \textit{Upapatti} --- Mathematical demonstrations and proofs
  \item \textit{Udāharaṇa} --- Numerical examples and worked problems
\end{enumerate}
\end{center}

\vspace{1cm}
\begin{center}
\longline

\textbf{\Large\dev{॥ इति समाप्तम् ॥}}

\vspace{0.5cm}
\textit{End of Part 1, Chunk 6 --- Pages 701--838}\\[0.3cm]
\textit{Bilingual Sanskrit-English Edition}

\vspace{0.3cm}
\longline

\vspace{0.5cm}
\begin{paracol}{2}
\begin{leftcolumn*}
\begin{center}
\dev{श्रीब्रह्मगुप्ताचार्येण विरचितः}\\
\dev{ब्राह्मस्फुटसिद्धान्तः}\\[0.3cm]
\dev{प्रथमभागस्य षष्ठः संचिकाङ्गं समाप्तम्}\\
\dev{(पृष्ठानि ७०१--८३८)}
\end{center}
\end{leftcolumn*}
\begin{rightcolumn}
\begin{center}
\textit{The Brāhmasphuṭasiddhānta}\\
\textit{composed by Śrī Brahmaguptācārya}\\[0.3cm]
\textit{The sixth portion of the first part is concluded}\\
\textit{(Pages 701--838)}
\end{center}
\end{rightcolumn}
\end{paracol}

\vspace{0.5cm}
\longline
\end{center}

\end{document}
%% ========================================================================
%% END OF DOCUMENT
%% =================================================